\begin{table}[htbp]
\centering
\footnotesize
\caption{Scomposizione per tipo di disposizione: quali clausole avrebbero un meccanismo commerciale?}
\label{tab:subindices}
\begin{threeparttable}
\begin{tabular}{lcc}
\toprule
Sotto-indice & $\times$ Verde & $\times$ Sporco \\
\midrule
Liberalizzazione dei beni verdi (WB) & -0.0115 & -0.0876$^{*}$ \\
 & (0.0879) & (0.0483) \\
 & [0.896] & [0.071] \\
\addlinespace
Accesso al mercato per i beni verdi (TREND) & -0.0271 & -0.0491$^{**}$ \\
 & (0.0303) & (0.0243) \\
 & [0.372] & [0.045] \\
\addlinespace
Meccanismo di risoluzione controversie (WB) & -0.0048 & 0.0067 \\
 & (0.0394) & (0.0525) \\
 & [0.903] & [0.898] \\
\addlinespace
Meccanismo di risoluzione controversie (TREND) & 0.0042 & -0.0052 \\
 & (0.0148) & (0.0140) \\
 & [0.778] & [0.709] \\
\addlinespace
Disposizioni vincolanti (TREND) & 0.0022 & -0.0013 \\
 & (0.0047) & (0.0037) \\
 & [0.638] & [0.724] \\
\addlinespace
Disposizioni di sola cooperazione (TREND) & 0.0132 & 0.0019 \\
 & (0.0121) & (0.0108) \\
 & [0.275] & [0.858] \\
\addlinespace
Clausole di spazio regolatorio (TREND) & 0.0242$^{**}$ & 0.0225$^{***}$ \\
 & (0.0099) & (0.0086) \\
 & [0.015] & [0.010] \\
\addlinespace
\midrule
Celle (tutte le righe) & \multicolumn{2}{c}{3,681,023} \\
\bottomrule
\end{tabular}
\begin{tablenotes}[flushleft]\footnotesize
\item Variante baseline, pannello collassato. Ogni riga \`e una stima separata in cui l'indice complessivo \`e sostituito dal sotto-indice indicato; il controllo di profondit\`a resta incluso.
\item \textbf{L'obiezione a cui risponde.} Un indice che somma tutte le disposizioni mescola clausole che potrebbero davvero spostare il commercio (abbassare i dazi sui beni verdi, imporre standard vincolanti) con dichiarazioni di intenti che non hanno alcun meccanismo. Se l'effetto esiste, dovrebbe concentrarsi nelle prime.
\item \textbf{Il limite, che \`e esso stesso un risultato.} Negli accordi cinesi del periodo le disposizioni dotate di un meccanismo commerciale sono rarissime: i due sotto-indici della Banca Mondiale che ne hanno uno sono diversi da zero in soli tre anni-paese e risultano perfettamente sovrapposti fra loro. Non \`e quindi possibile distinguerne gli effetti: non per un difetto del metodo, ma perch\'e quel contenuto negli accordi quasi non c'\`e.
\item Le righe \emph{sola cooperazione} e \emph{spazio regolatorio} funzionano da controllo negativo: sono clausole senza meccanismo commerciale, e l\`i non dovrebbe emergere nulla.
\item Errore standard fra parentesi tonde, $p$-value fra parentesi quadre. $^{*}$ $p<0.10$, $^{**}$ $p<0.05$, $^{***}$ $p<0.01$.
\end{tablenotes}
\end{threeparttable}
\end{table}
